\chapter{Временные ряды и последовательный анализ}
\label{ch:timeseries}

Главы~\ref{ch:basics}--\ref{ch:regression} опирались на предпосылку независимости (или, в регрессии, на некоррелированность остатков после учёта предикторов). Временной ряд~--- последовательность наблюдений, упорядоченных во времени. Зависимость между соседними отсчётами нарушает i.i.d.-модель и требует отдельного аппарата.
Прогнозирование требует явного учёта тренда, сезонности и автокорреляции. Простая регрессия на время $t$ часто оставляет структуру в остатках, и $t$-критерии для коэффициентов теряют обоснованность.

Последовательный анализ (SPRT) решает иную, но родственную задачу: как проверять гипотезы по мере поступления данных, не фиксируя объём выборки заранее. Это критично в A/B-тестах и мониторинге процессов, где подглядывание в накопленные результаты без поправки раздувает уровень ошибки I рода~--- та же частотная логика, что в гл.~\ref{ch:multiple}, но в динамической постановке.
Настоящая глава связывает классическое прогнозирование (декомпозиция, сглаживание, ARIMA) с последовательными процедурами Вальда и доверительными интервалами Стейна.
Читатель, освоивший регрессию, уже знаком с идеей остатков без структуры. Во временных рядах эта идея переносится на последовательные ошибки прогноза, а нарушение независимости~--- на автокорреляцию.

\section{Постановка прогнозирования}

Пусть $y_1,\ldots,y_T$~--- ряд с постоянным шагом между отсчётами, цель~--- прогноз $\hat y_{T+d}$ на горизонт $d\in\{1,\ldots,D\}$.
Ошибка прогноза $e_t=y_t-\hat y_{t\mid t-1}$. Хороший одношаговый прогноз даёт остатки, близкие к белому шуму~--- аналог хороших остатков в регрессии, но теперь проверяем отсутствие автокорреляции, а не только нулевое среднее.
Регрессия только на время $t$ часто оставляет структуру в остатках. Модели с лагами $y_{t-1},\ldots$ и сезонными компонентами учитывают автокорреляцию явно.
Критерий качества~--- не $R^2$ на обучении, а ошибка на отложенном хвосте или скользящая валидация. Переобучение на коротком ряду даёт оптимистичные прогнозы, сходное с переобучением в регрессии при $n\approx k$.
Разделение на обучающий и тестовый отрезки во времени должно уважать хронологию. Будущее не может попасть в обучение~--- иначе оценка точности прогноза будет завышена.

Наивные методы $\hat y_{T+d}=y_T$, $\hat y_{T+d}=\bar y$, $\hat y_{T+d}=y_{T+d-s}$ (сезонный наивный, период~$s$).
Наивные прогнозы~--- baseline. Любая сложная модель должна превосходить их на тестовом отрезке, иначе усложнение не оправдано.
Регрессия на лаги $y_t=\beta_0+\sum_{j=1}^p\beta_j y_{t-j}+\varepsilon_t$.
Выборочная ACF $r_\tau=\dfrac{\sum_{t=1}^{T-\tau}(y_t-\bar y)(y_{t+\tau}-\bar y)} {\sum_{t=1}^T(y_t-\bar y)^2}$.
ACF измеряет линейную связь ряда с его сдвигами~--- обобщение корреляции Пирсона (гл.~\ref{ch:anova}) на упорядоченные пары $(y_t,y_{t+\tau})$.
В отличие от корреляции в регрессии, здесь переменные~--- один и тот же процесс на разных лагах. Значимые $r_\tau$ указывают, что прошлое несёт информацию о настоящем.
Частная автокорреляция (PACF) отвечает на вопрос, какова связь $y_t$ и $y_{t-\tau}$ после учёта промежуточных лагов~--- ключ к идентификации порядка AR.

Пример: продажи вина. Ряд продаж австралийского вина содержит ежемесячные продажи вина (1980--1994). Видны тренд и годовая сезонность ($s=12$).
Автокорреляция на лагах 1, 12 и~24 отражает инерцию и годовой цикл (рис.~\ref{fig:ts-acf}).
Пики на лагах, кратных 12,~--- типичный отпечаток сезонности. Медленное затухание ACF на малых лагах указывает на нестационарность или длинную память, что направляет выбор дифференцирования и SARIMA.
Сравнение с корреляцией в регрессии (гл.~\ref{ch:anova}): там $r$ измеряет связь двух разных переменных по одному индексу $i$. Здесь~--- связь одного процесса с самим собой во времени, и нарушение независимости по $t$~--- норма, а не исключение.
Типовой маршрут для продаж вина: график ряда $\to$ декомпозиция $\to$ ADF/KPSS $\to$ SARIMA или Хольт--Уинтерс $\to$ прогноз на отложенном хвосте $\to$ Q Льюнга--Бокса на остатках. Если после подбора модели $Q$ значим, высокий на обучении $R^2$ не спасает. В остатках осталась структура, и прогноз на $d>1$ будет систематически смещён.

\begin{figure}[htbp]
    \centering
    \includegraphics[width=\FigWidth]{\FigDir/chapter-6/wineacf.png}
    \caption{Автокорреляционная функция (ACF) ряда.}
    \label{fig:ts-acf}
\end{figure}

\section{Декомпозиция и стационарность}

Ряд часто представляют как сумму или произведение компонент: тренд $T_t$, сезонность $S_t$, цикл $C_t$, остаток $R_t$.
Декомпозиция~--- содержательный шаг. Отделить медленное (тренд, цикл) от быстрого (сезонность, шум) и понять, какая часть ряда подлежит прогнозированию, а какая~--- моделированию как случайный процесс.

Аддитивная декомпозиция задаёт $y_t=T_t+S_t+C_t+R_t$.
Уместна, когда амплитуда сезонных колебаний не зависит от уровня ряда.
Мультипликативная декомпозиция.~--- $y_t=T_t\cdot S_t\cdot C_t\cdot R_t$ (или $y_t=T_t S_t R_t$ без отдельного цикла).
После логарифмирования мультипликативная модель становится аддитивной: $\ln y_t=\ln T_t+\ln S_t+\ln R_t$.

STL-декомпозиция оценивает $T_t$, $S_t$, $R_t$ итеративно, допуская медленное изменение сезонного профиля.
Снятие сезонности: $y_t^{\mathrm{deseas}}=y_t-S_t$ (аддитив) или $y_t/S_t$ (мультипликатив) упрощает анализ тренда.
Преобразование Бокса--Кокса (гл.~\ref{ch:regression}) часто применяют до декомпозиции, если дисперсия растёт с уровнем ряда.

\begin{definition}
Ряд $\{y_t\}$ называется слабо стационарным, если $\mathbb{E}y_t$ и $\cov(y_t,y_{t-k})$ не зависят от~$t$ для всех $k$.
\end{definition}

Ряды с трендом, сезонностью или меняющейся дисперсией нестационарны.
Стационарность~--- аналог одинакового распределения для каждого момента времени. Среднее и автоковариации не плывут.
ARMA-модели строят для стационарного ряда. Нестационарность устраняют дифференцированием.
Практически ряд редко бывает строго стационарным на длинном горизонте (смена рынка, технологий), но локальная стационарность после дифференцирования~--- рабочая аппроксимация для прогноза на коротком горизонте.
Практически стационарность проверяют не одним тестом. ADF с $H_0$ о единичном корне и KPSS с $H_0$ о стационарности дополняют друг друга. Расхождение результатов~--- повод посмотреть на график дифференцированного ряда. Возможно, достаточно $\nabla_{12}$ без лишнего $\nabla$, или наоборот, сезонный тренд требует $D=1$.
Избыточное дифференцирование порождает лишнюю автокорреляцию и ухудшает долгосрочный прогноз.

Стабилизация дисперсии. Преобразование Бокса--Кокса $y^{(\lambda)}=\frac{y^\lambda-1}{\lambda}$ при $\lambda\neq 0$, $\log y$ при $\lambda=0$.
Устранение тренда $\nabla y_t=y_t-y_{t-1}$, сезонное $\nabla_s y_t=y_t-y_{t-s}$.
Чем меньше порядок дифференцирования, тем меньше дисперсия долгосрочного прогноза. Для сезонных рядов начинают с $\nabla_s$.
KPSS ($H_0\colon$ стационарность): \[ \mathrm{KPSS}=\frac{1}{T^2}\sum_{i=1}^T\Bigl(\sum_{t=1}^i\varepsilon_t\Bigr)^2\Big/\hat\lambda^2, \] $\varepsilon_t$~--- остатки тестовой регрессии, $\hat\lambda^2$~--- оценка спектра на нулевой частоте.
Отклонение $H_0$~--- признак нестационарности.
ADF и модификации (Дики--Фуллера, Филлипса--Перрона) проверяют $H_0\colon$ наличие единичного корня (нестационарность).
KPSS и ADF дополняют друг друга. Важно, что именно является $H_0$.
При расхождении результатов смотрят на график ряда и ACF. Формальные тесты~--- ориентир, а не единственный аргумент. Избыточное дифференцирование ($d$ или $D$ слишком велики) вводит лишнюю автокорреляцию в дифференцированный ряд и ухудшает прогноз. Поэтому предпочитают минимальный порядок, при котором KPSS не отвергает стационарность (или ADF отвергает единичный корень).
Для ряда меток времени событий перед моделированием агрегируют counts по интервалам. Неравномерная сетка требует специальных моделей (point process), здесь~--- иллюстрация подготовки данных.

\section{Экспоненциальное сглаживание}

Простое сглаживание (SES). Прогноз одного шага вперёд: \[ \hat y_{t+1\mid t}=\alpha y_t+(1-\alpha)\hat y_{t\mid t-1}, \quad \alpha\in(0,1].
\] Эквивалентно $\hat y_{t+1\mid t}=l_t$, где $l_t=\alpha y_t+(1-\alpha)l_{t-1}$.
Подходит для рядов без тренда и сезонности. Горизонт $d>1$~--- плоский прогноз $\hat y_{t+d\mid t}=l_t$.
Параметр $\alpha$ задаёт память модели. Большие $\alpha$~--- быстрая реакция на новые данные, малые~--- инерция.
SES можно рассматривать как простейший прогноз без явной модели автокорреляции. Экспоненциально убывающие веса прошлых наблюдений.

Метод Хольта (тренд).\[ \hat y_{t+d\mid t}=l_t+d\,b_t, \quad l_t=\alpha y_t+(1-\alpha)(l_{t-1}+b_{t-1}), \quad b_t=\beta(l_t-l_{t-1})+(1-\beta)b_{t-1}.
\] Мультипликативный тренд: $\hat y_{t+d\mid t}=l_t b_t^d$.
Компонента $b_t$~--- сглаженная скорость изменения уровня. При экстраполяции на большой $d$ линейный или мультипликативный тренд может давать нереалистичные прогнозы, если структура ряда сменилась.

Хольт--Уинтерс (сезонность). Аддитивная сезонность периода $m$: \[ \hat y_{t+d\mid t}=l_t+d\,b_t+s_{t-m+(d\bmod m)}, \] \[ l_t=\alpha(y_t-s_{t-m})+(1-\alpha)(l_{t-1}+b_{t-1}),\quad b_t=\beta(l_t-l_{t-1})+(1-\beta)b_{t-1},\quad s_t=\gamma(y_t-l_{t-1}-b_{t-1})+(1-\gamma)s_{t-m}.
\] Мультипликативная: $\hat y_{t+d\mid t}=(l_t+d b_t)\,s_{t-m+(d\bmod m)}$, $l_t=\alpha y_t/s_{t-m}+(1-\alpha)(l_{t-1}+b_{t-1})$, $s_t=\gamma y_t/(l_{t-1}+b_{t-1})+(1-\gamma)s_{t-m}$.
Для продаж вина сезонный метод даёт прогноз, повторяющий годовой профиль.

\begin{figure}[htbp]
    \centering
    \includegraphics[width=\FigWidth]{\FigDir/chapter-6/fig_7_from_statsmodels.png}
    \caption{Наблюдения и прогноз на конце ряда.}
    \label{fig:ts-forecast}
\end{figure}

Параметры $\alpha,\beta,\gamma$ подбирают по минимуму SSE на обучающем отрезке. Переобучение на коротком ряду завышает точность на обучении.
Экстраполяция тренда за пределы истории ненадёжна.
Мультипликативная сезонность требует $y_t>0$. При нулях используют аддитивную форму или сдвиг ряда.
Сглаживание удобно для оперативных прогнозов и презентации тренда. ARIMA~--- когда нужна явная модель автокорреляции и статистические интервалы с обоснованием. На практике часто сравнивают оба подхода на отложенном хвосте и выбирают по RMSE или MAPE.
Сравнение SES, Хольта--Уинтерса и SARIMA на одном ряде~--- типовой пример выбора модели. Наивный прогноз задаёт baseline. Если Хольт--Уинтерс не бьёт последнее значение на тесте, сезонность или тренд моделируются неверно. Содержательный критерий~--- не минимум AIC на обучении, а ошибка на последних $h$ точках, которые модель не видела при подборе.

\section{ARIMA-модели}

Класс ARIMA \citep{box2015} обобщает авторегрессию и скользящее среднее.

Оператор запаздывания $B$: $By_t=y_{t-1}$, $\nabla=1-B$.

Стационарный AR($p$): \[ \phi(B)y_t=(1-\phi_1 B-\cdots-\phi_p B^p)y_t=\varepsilon_t, \quad \varepsilon_t\sim WN(0,\sigma_\varepsilon^2).
\] MA($q$): $y_t=\theta(B)\varepsilon_t=\varepsilon_t+\theta_1\varepsilon_{t-1}+\cdots+\theta_q\varepsilon_{t-q}$.
ARMA ($p,q$): \[ y_t=\phi_1 y_{t-1}+\cdots+\phi_p y_{t-p} +\varepsilon_t+\theta_1\varepsilon_{t-1}+\cdots+\theta_q\varepsilon_{t-q}, \] или $\phi(B)y_t=\theta(B)\varepsilon_t$.
Теорема Вольда: стационарный ряд аппроксимируется ARMA с любой точностью.
White noise $\varepsilon_t$~--- аналог хороших остатков регрессии. Нулевое среднее, постоянная дисперсия, отсутствие автокорреляции.
Контраст с регрессией: там хорошие остатки~--- i.i.d.
$N(0,\sigma^2)$ по $i$. Здесь~--- белый шум по $t$. Нарушение независимости по времени (автокорреляция остатков прогноза) означает, что модель не выжала всю предсказуемую информацию из прошлого ряда.
Условие стационарности AR($p$): корни характеристического многочлена $\phi(z)=0$ лежат вне единичного круга. Для MA($q$)~--- обратимость.
Нарушение стационарности означает, что дисперсия $y_t$ не ограничена или ряд взрывается~--- признак необходимости дифференцирования.

\begin{definition}
Модель ARIMA$(p,d,q)$ задаёт, что $\nabla^d y_t$ удовлетворяет ARMA$(p,q)$, то есть $\phi(B)(1-B)^d y_t=\theta(B)\varepsilon_t$, где $\varepsilon_t$~--- белый шум.
\end{definition}

SARIMA ($p,d,q$)$\times$($P,D,Q$)$_s$ добавляет сезонные операторы периода $s$: \[ \phi(B)\Phi(B^s)(1-B)^d(1-B^s)^D y_t =\theta(B)\Theta(B^s)\varepsilon_t, \] где $\Phi(B^s)=1-\Phi_1 B^s-\cdots-\Phi_P B^{Ps}$, $\Theta(B^s)=1+\Theta_1 B^s+\cdots+\Theta_Q B^{Qs}$.
Для месячных продаж вина типичны $s=12$, $D=1$ (сезонное дифференцирование).
Пара $(p,d,q)$ и $(P,D,Q)_s$ подбирают итеративно. Графики ACF/PACF дифференцированного ряда, затем минимизация AIC по сетке.
Переусложнение (слишком большие $p,q$) даёт нестабильные прогнозы на длинном горизонте, хотя на обучении подгонка улучшается.

ACF/PACF для идентификации. AR($p$)~--- PACF обрывается после лага $p$. MA($q$)~--- ACF обрывается после $q$.
ARIMA($p,d,0$): PACF значима до лага $p$. ARIMA($0,d,q$): ACF до $q$.
Подбор $p,q,P,Q$. Сетка по AIC/AICc. Стартовые $p,q$~--- последние значимые лаги ACF/PACF.
Прогноз ARIMA. В уравнении заменить $t$ на $T+h$, будущие $y$~--- на прогнозы, будущие $\varepsilon$~--- на 0, прошлые $\varepsilon$~--- на остатки. Минимизирует MSE при правильной спецификации.
Интервалы прогноза. При гауссовых $\varepsilon$ дисперсия ошибки на горизонте $h$ растёт. Интервалы $\hat y_{T+h\mid T}\pm z_{1-\alpha/2}\,\hat\sigma_h$.
После преобразования Бокса--Кокса точечный прогноз обратно преобразуют. Интервалы сильнее зависят от $\lambda$.
Диагностика. Q-критерий Льюнга--Бокса на остатках прогноза.

Критерий Льюнга--Бокса. Для остатков прогноза $\varepsilon_1,\ldots,\varepsilon_T$ проверяют $H_0\colon r_1=\cdots=r_L=0$: \[ Q=T(T+2)\sum_{\tau=1}^L\frac{r_\tau^2}{T-\tau}\sim\chi^2_{L-K}, \] $K$~--- число оценённых параметров модели.
Большое $Q$ указывает на неучтённую автокорреляцию (модель можно улучшить), но не гарантирует, что улучшение достижимо в выбранном классе.
Вычитание $K$ из степеней свободы~--- поправка на оценённые параметры, аналогичная использованию $n-k-1$ в регрессии вместо $n$.
Если $Q$ значим при разумном $L$, пробуют увеличить порядок AR или MA, добавить сезонные компоненты или внешние регрессоры. Если $Q$ мал, но прогноз на тесте плох, проблема может быть в структурном сдвиге, а не в порядке модели.

ARIMA плохо переносит структурные сдвиги (COVID, смена политики).
Календарные эффекты (праздники, число рабочих дней) добавляют регрессоры в рамках ARIMAX.
Остатки прогноза должны быть несмещёнными, стационарными и неавтокоррелированными. Нормальность желательна для интервалов, но не обязательна для точечного прогноза.
Сравнение с регрессией: там хорошие остатки~--- i.i.d. $N(0,\sigma^2)$. Здесь~--- белый шум $\varepsilon_t$ без автокорреляции. Нормальность $\varepsilon_t$ нужна для симметричных интервалов прогноза, но точечный прогноз $\hat y_{T+h\mid T}$ опирается на условное матожидание и состоятелен при более широких предпосылках.
Идентификация $(p,d,q)$ по ACF/PACF~--- эвристика, а не теорема. На коротком ряде пики могут быть шумом. Типовой приём: подобрать 2--3 кандидата по AIC, сравнить прогноз на отложенной выборке, затем Q-тест на остатках победителя. Если два кандидата близки по RMSE, предпочитают более простую модель (меньше $(p,q)$)~--- та же философия, что штраф $k$ в AIC/BIC (гл.~\ref{ch:regression}).

\section{Последовательный анализ (SPRT)}

Классический тест фиксирует $n$ заранее. Последовательный анализ Вальда накапливает данные до принятия решения, контролируя $\alpha$ и $\beta$.
В A/B-тестировании это позволяет остановить эксперимент, как только накоплено достаточно доказательств без подглядывания в фиксированный $t$-тест~--- ключевое отличие от наивного мониторинга p-value.
Связь с гл.~\ref{ch:multiple}: многократная проверка гипотез без поправки раздувает FWER. SPRT встроенно учитывает последовательность решений, если границы заданы до начала эксперимента.

SPRT для доли Бернулли $X_i\sim Ber(p)$ i.i.d., коридор $p_L\le p_0\le p_U$.
\[ H_0\colon p\le p_L, \qquad H_1\colon p\ge p_U.
\] Статистика $d_m=\sum_{i=1}^m X_i$, границы \[ A=\frac{1-\beta}{\alpha}, \quad B=\frac{\beta}{1-\alpha}, \] \[ a_m=\frac{\ln B+m\ln\frac{1-p_L}{1-p_U}}{\ln\frac{p_U}{p_L}-\ln\frac{1-p_U}{1-p_L}}, \quad r_m=\frac{\ln A+m\ln\frac{1-p_L}{1-p_U}}{\ln\frac{p_U}{p_L}-\ln\frac{1-p_U}{1-p_L}}.
\] Если $d_m\ge r_m$~--- отвергаем $H_0$. Если $d_m\le a_m$~--- принимаем $H_0$. Иначе $m\leftarrow m+1$.
Момент остановки $N$~--- случайная величина.
Ожидаемый объём выборки (оперативная характеристика $L(p)$): \[ \mathbb{E}_p(N)=\frac{L(p)\ln B+(1-L(p))\ln A} {p\ln\frac{p_U}{p_L}+(1-p)\ln\frac{1-p_U}{1-p_L}}, \] $L(p)=\dfrac{A^h-1}{A^h-B^h}$, $h$~--- корень уравнения для $p$.
При $p=p_L$ или $p=p_U$ средний $N$ существенно меньше фиксированного $n$ при тех же $\alpha,\beta$.
Коридор $[p_L,p_U]$ задаёт практически значимый эффект, а не точечное $p_0$~--- та же философия, что при выборе MDE в планировании эксперимента.
Отношение правдоподобий накапливается по мере поступления $X_i$. Границы $a_m,r_m$~--- логарифмические пороги для этого отношения, аналогичные критическим значениям $t$- или $F$-критерия в фиксированной выборке.

Усечение и группировка (batching). Если при $m=n_0$ решение не принято, применяют усечение. Сравнивают $d_m$ с серединой $(a_{n_0}+r_{n_0})/2$.
Это нарушает точный контроль $\alpha,\beta$, но ограничивает максимальный $n$.
В промышленном мониторинге усечение~--- компромисс между строгим SPRT и жёстким дедлайном по времени или объёму данных.

Группировка. Наблюдения поступают пачками по $v$ штук. Сравнение с $a_m,r_m$ только при $m=v,2v,\ldots$.
Увеличивает типичный $N$. Истинные ошибки могут превышать номинальные, но $\alpha'\le\alpha/(1-\beta)$, $\beta'\le\beta/(1-\alpha)$, что при малых $\alpha,\beta$ близко к номиналу.

SPRT для нормального среднего. При известной $\sigma^2$ используют $d_m=\sum X_i$ с границами \[ a_m=\frac{\sigma^2}{\mu_U-\mu_L}\ln B+m\frac{\mu_U+\mu_L}{2},\quad r_m=\frac{\sigma^2}{\mu_U-\mu_L}\ln A+m\frac{\mu_U+\mu_L}{2} \] для $H_0\colon\mu\le\mu_L$ против $H_1\colon\mu\ge\mu_U$.
При неизвестной $\sigma$~--- последовательные $t$-процедуры, группировка или подстановка пилотной оценки $\sigma$.

A/B-тестирование. Конверсия~--- SPRT Бернулли. Средний чек при нормальном приближении~--- SPRT для среднего.
Перед запуском фиксируют $p_L,p_U$ (практически значимый эффект), а не точечное $p_0$.
Типичные ошибки. Многократное подглядывание в фиксированный $t$-тест без поправки раздувает $\alpha$. Слишком узкий коридор увеличивает $\mathbb{E}N$. Игнорирование сезонности в метриках. A/A-тест для проверки инфраструктуры. Разные метрики (CTR, ARPU) требуют разных моделей. Нарушение порядка случайизации (SUTVA) ломает интерпретацию.
Последовательный мониторинг с SPRT или последовательными ДИ сохраняет частотный контроль при корректной постановке.
В отличие от фиксированного $n$-теста, где мощность известна до сбора данных, в SPRT момент остановки случаен. Отчёт должен указывать не только решение ($H_0$ или $H_1$), но и фактическое $N$, иначе воспроизводимость страдает. Типовой сценарий A/B-теста: фиксируем $p_L,p_U$ (MDE) до запуска, считаем границы SPRT, останавливаемся при пересечении~--- это альтернатива ежедневному мониторингу p-value $t$-теста без поправки.
Сравнение с гл.~\ref{ch:multiple}: оба подхода борются с многократной проверкой, но SPRT встроен в дизайн, а апостериорные поправки~--- в анализ.

Последовательные доверительные интервалы. Проблема фиксированного $n$. При неизвестной $\sigma^2$ нельзя построить интервал для $\mu$ фиксированной длины $2d$ с покрытием $1-\alpha$ для всех $\mu,\sigma$ (Данциг, 1940).

Двухэтапная процедура Стейна. Пилот $X_1,\ldots,X_m$, $m\ge 2$, $S_m^2$~--- выборочная дисперсия.
\[ \hat C=\frac{t_{1-\alpha/2,m-1}^2 S_m^2}{d^2},\quad N=\max(\lfloor\hat C\rfloor+1,m),\quad J_N=[\bar X_N-d,\bar X_N+d].
\] $P(\mu\in J_N)\ge 1-\alpha$. При $d\to 0$ покрытие $\to 1-\alpha$, но $\mathbb{E}(N)/C\to t_{1-\alpha/2,m-1}^2/z_{1-\alpha/2}^2>1$ (неэффективно).
Двухэтапная процедура разделяет оценку дисперсии и основную выборку. Полностью последовательная~--- экономит наблюдения, останавливаясь, как только накопленная информация достаточна для интервала длины $2d$.

Полностью последовательная процедура. Для $n=m,m+1,\ldots$ $\hat C_n=z_{1-\alpha/2}^2 S_n^2/d^2$. Остановка при $n\ge\hat C_n$.
Асимптотически эффективна: $\mathbb{E}(N)\le C+m$, $\mathbb{E}(N)/C\to 1$.

Разность двух средних ($\sigma_1=\sigma_2=\sigma$): пилотные выборки объёма $m$, $U_m^2=(S_{1m}^2+S_{2m}^2)/2$, \[ \hat C=\frac{2t_{1-\alpha/2,2m-2}^2 U_m^2}{d^2},\quad N=\max(\lfloor\hat C\rfloor+1,m), \] интервал $J_N=[T_N-d,T_N+d]$, $T_N=\bar X_{1N}-\bar X_{2N}$.
При $\sigma_1\neq\sigma_2$ используют квантиль $h_{m,1-\alpha/2}$ распределения разности двух $t$ и отдельные $\hat C_1,\hat C_2$ для объёмов двух групп (решение задачи Беренса--Фишера в последовательной постановке).

Пример: ADHD и реакция. Временной ряд активности (лабораторный пример) содержит измерения времени реакции. Для демонстрации SPRT можно тестировать, улучшает ли вмешательство среднее время относительно контроля, задавая $\mu_L,\mu_U$ как практически значимые границы эффекта, а не точечное $\mu_0$.
Связь с временными рядами здесь двойная. Данные упорядочены во времени, но SPRT предполагает i.i.d. внутри накопленной выборки~--- автокорреляцию остатков нужно учесть отдельно или агрегировать наблюдения в независимые блоки.
Если метрика A/B-теста~--- дневная конверсия с сильной недельной сезонностью, наивный SPRT по дням может увидеть ложный эффект. Агрегация по неделям или включение сезонных регрессоров~--- содержательная часть постановки.

\begin{itemize}
    \item Подглядывание (peeking) в классический t-тест без поправки~---
          ложные срабатывания.
    \item SPRT требует заранее заданных $p_L,p_U$ (или $\mu_L,\mu_U$). Слишком узкий
          коридор увеличивает $\mathbb{E}N$.
    \item Группировка уменьшает число взглядов, но снижает мощность.
    \item Сезонность в продаж вина игнорируется простым AR(1)~---
          систематические ошибки прогноза.
    \item Останавливать A/B-тест, когда уже видно победителя, без SPRT~---
          завышенная доля ложных побед.
\end{itemize}
Содержательный итог главы: временной ряд и SPRT решают одну проблему~--- данные приходят по порядку, и порядок нельзя игнорировать.
В прогнозе порядок задаёт автокорреляцию. В A/B-тесте~--- многократные взгляды на накопленную статистику. В обоих случаях формальная процедура (ARIMA-диагностика, SPRT) заменяет наивную экстраполяцию или наивный мониторинг p-value.

Пример: продажи и A/B-тест. Ежемесячные продажи ($T=60$) с трендом и $s=12$. STL выделяет сезонный профиль. ADF на исходном ряду не отвергает единичный корень, после $\nabla_{12}$~--- KPSS не отвергает стационарность. SARIMA$(1,0,1)\times(1,1,1)_{12}$ и Хольт--Уинтерс сравнивают на последних 12 точках. RMSE сопоставим, но SARIMA даёт более узкие интервалы прогноза. Q Льюнга--Бокса на остатках SARIMA не значим.

Параллельно продуктовая команда ведёт A/B-тест конверсии ($p_L=0{,}02$, $p_U=0{,}05$ по SPRT Бернулли). На 15-й день накоплено $d_{15}=3$ успехов из 120 визитов в варианте B против 1 из 118 в A. Границы SPRT ещё не пересечены~--- решение подождать корректно, в отличие от ежедневного подглядывания в p-value $z$-теста на разность двух долей.

Связь разделов: и ARIMA, и SPRT требуют заранее зафиксированной процедуры и уважения к порядку данных. В отчёте по ряду указывают RMSE на отложенной выборке и Q. По A/B~--- $p_L,p_U$, фактическое $N$ остановки и MDE. Обе линии главы~--- про контроль ложных выводов при последовательном поступлении информации.

\section{Расширенная диагностика рядов и SPRT}

Чеклист диагностики остатков прогноза. Остатки прогноза $\hat\varepsilon_t=y_t-\hat y_{t\mid t-1}$ должны вести себя как белый шум. После подбора ARIMA/SARIMA или ETS пройдите единый контрольный перечень:

\begin{enumerate}
  \item Несмещённость (необходимо): $\bar{\hat\varepsilon}\approx 0$.
        Формально~--- $t$-критерий или критерий Уилкоксона для среднего остатков.
  \item Стационарность (необходимо): график остатков без тренда
        и сезонной структуры. KPSS с $H_0\colon$ стационарность: \[ \mathrm{KPSS}=\frac{1}{T^2}\sum_{i=1}^T\Bigl(\sum_{t=1}^i\hat\varepsilon_t\Bigr)^2\Big/\hat\lambda^2, \] $\hat\lambda^2$~--- оценка спектра на нулевой частоте. Отклонение $H_0$~--- нестационарные остатки.
  \item Неавтокоррелированность (необходимо): коррелограмма остатков.
        Q Льюнга--Бокса для $H_0\colon r_1=\cdots=r_L=0$.
  \item Нормальность (желательно для симметричных интервалов):
        Q--Q график. Шапиро--Уилка на малых $T$.
\end{enumerate}

\begin{critcard}{Q Льюнга--Бокса (остатки прогноза)}
Ряд & $\hat\varepsilon_1,\ldots,\hat\varepsilon_T$ \\
$H_0$ & $r_1=\cdots=r_L=0$ (нет автокорреляции до лага $L$) \\
Статистика & $Q=T(T+2)\sum_{\tau=1}^L \dfrac{r_\tau^2}{T-\tau}$ \\
Нулевое распр. & $\chi^2_{L-K}$, $K$~--- число оценённых параметров модели \\
Интерпретация & большой $Q$~--- в остатках осталась предсказуемая структура
\end{critcard}

\begin{critcard}{KPSS (стационарность остатков)}
Ряд & $\hat\varepsilon_1,\ldots,\hat\varepsilon_T$ \\
$H_0$ & стационарность (в отличие от ADF, где $H_0$~--- единичный корень) \\
Статистика & $\mathrm{KPSS}=\dfrac{1}{T^2}\sum_{i=1}^T\bigl(\sum_{t=1}^i\hat\varepsilon_t\bigr)^2\big/\hat\lambda^2$ \\
Решение & отклонение $H_0$~--- тренд/сезонность в остатках. Доработать модель \\
Замечание & KPSS и ADF дополняют друг друга. Следите, что является $H_0$
\end{critcard}

\begin{critcard}{Шапиро--Уилка (нормальность остатков, малые $T$)}
Выборка & $\hat\varepsilon_1,\ldots,\hat\varepsilon_T$ \\
$H_0$ & остатки нормальны \\
Статистика & $W$ (корреляция с ожидаемыми порядковыми статистиками) \\
Интерпретация & малый p-value~--- тяжёлые хвосты. Точечный прогноз может оставаться приемлемым, интервалы~--- нет
\end{critcard}

Порядок на практике: график остатков $\to$ KPSS $\to$ Q по лагам $1\ldots L$ $\to$ Q--Q и Шапиро. Значимый Q при хорошем RMSE на отложенной выборке часто означает, что класс моделей исчерпан, а не что подгонка плохая.

Двусторонний SPRT для среднего нормального ряда. При известной $\sigma^2$ и симметричном коридоре вокруг $\mu_0$: $\bigl|(\mu-\mu_0)/\sigma\bigr|\leq\delta$ против $\bigl|(\mu-\mu_0)/\sigma\bigr|>\delta$.

\begin{critcard}{Двусторонний SPRT, $X_i\sim N(\mu,\sigma^2)$, $\sigma$ известна}
Выборка & $X_1,\ldots,X_m$ по мере поступления \\
$H_0$ & $\bigl|(\mu-\mu_0)/\sigma\bigr|\leq\delta$ \\
$H_1$ & $\bigl|(\mu-\mu_0)/\sigma\bigr|>\delta$ \\
Статистика & $d_m=\ln\cosh\!\left(\dfrac{\delta}{\sigma}\sum_{i=1}^m(X_i-\mu_0)\right)$ \\
Границы & $a_m=\ln B+m\dfrac{\delta^2}{2}$, \quad $r_m=\ln A+m\dfrac{\delta^2}{2}$ \\
Правило & $d_m\geq r_m$~--- $H_1$. $d_m\leq a_m$~--- $H_0$. Иначе $m:=m+1$ \\
Константы & $A=(1-\beta)/\alpha$, $B=\beta/(1-\alpha)$
\end{critcard}

Односторонний коридор $\mu_L\leq\mu_0\leq\mu_U$ (раздел SPRT в этой главе): $d_m=\sum_{i=1}^m X_i$, \[ a_m=\frac{\sigma^2}{\mu_U-\mu_L}\ln B+m\frac{\mu_U+\mu_L}{2},\qquad r_m=\frac{\sigma^2}{\mu_U-\mu_L}\ln A+m\frac{\mu_U+\mu_L}{2}.
\] При неизвестной $\sigma$~--- последовательный $t$-процесс или пилотная оценка $\sigma$.

Последовательный доверительный интервал фиксированной длины для медианы. При фиксированном $n$ интервал для медианы $\theta$ длины $2d$ с покрытием $1-\alpha$ для всех $F$ не существует (Данциг).
Полностью последовательная процедура. Для $X_i\sim F(x-\theta)$ с симметричным гладким $F$ набирают выборку, пока интервал по порядковым статистикам шире $2d$: \[ b(n)=\max\!\left(1,\left\lfloor\tfrac12(n-z_{1-\alpha/2}\sqrt{n}-1)\right\rfloor\right),\quad a(n)=n-b(n)+1, \] $J_n=[X_{(b(n))},\,X_{(a(n))}]$. Остановка при $X_{(a(n))}-X_{(b(n))}\leq 2d$.
Асимптотически $P(\theta\in J_N)\to 1-\alpha$ при $d\to 0$. $\mathbb{E}(N)$ сопоставим с $z_{1-\alpha/2}^2/(4d^2 f^2(0))$.

Для среднего при неизвестной $\sigma$ используют процедуры Стейна (двухэтапную или полностью последовательную с $\hat C_n=z_{1-\alpha/2}^2 S_n^2/d^2$)~--- см. раздел последовательных ДИ в основной части главы.

SPRT для дисперсии нормальной выборки. При известном $\mu$ проверяют коридор $\sigma_L\leq\sigma_0\leq\sigma_U$.

\begin{critcard}{SPRT для дисперсии, $\mu$ известно, $X_i\sim N(\mu,\sigma^2)$}
$H_0$ & $\sigma\leq\sigma_L$ \\
$H_1$ & $\sigma\geq\sigma_U$ \\
Статистика & $d_m=\sum_{i=1}^m (X_i-\mu)^2$ \\
Границы & $a_m=\dfrac{2\ln B+m\ln(\sigma_U^2/\sigma_L^2)}{1/\sigma_L^2-1/\sigma_U^2}$, \\
& $r_m=\dfrac{2\ln A+m\ln(\sigma_U^2/\sigma_L^2)}{1/\sigma_L^2-1/\sigma_U^2}$ \\
$\mu$ неизвестно & $d_m=\sum_{i=1}^m(X_i-\bar X)^2$. На шаге $m$ сравнивать с $a_{m-1},r_{m-1}$
\end{critcard}

Пример содержательной постановки: контроль, не превышает ли погрешность прибора заявленный уровень ($\sigma_0$~--- допуск производителя).

\section{Практика и семинар}
\label{sec:practice-ch6}

Практика семинаров 9--10: ARIMA для продаж вина и последовательный анализ в стиле A/B-теста.

Семинар~9: месячные продажи австралийского вина. Ряд \texttt{monthly-australian-wine-sales-th.csv} (1980--1994, $s=12$): тренд и выраженная годовая сезонность. Пики ACF на лагах $12,24,\ldots$.

Рекомендуемый маршрут (согласован с примером продаж вина в основной части главы):
\begin{enumerate}
  \item График ряда, STL-декомпозиция. При росте дисперсии с уровнем~---
        Бокса--Кокса или $\log y_t$.
  \item KPSS на исходном ряде (ожидаемо $H_0$ отвергается),
        сезонное дифференцирование $\nabla_{12}$, повторный KPSS.
  \item ACF/PACF дифференцированного ряда. Сетка SARIMA по AIC
        (напр.\ $(p,d,q)\times(P,D,Q)_{12}$).
  \item Сравнение с Хольтом--Уинтерс (аддитивный/мультипликативный
        сезон, $m=12$) на последних 12 точках отложенной выборки.
  \item Чеклист остатков победителя: KPSS, Q Льюнга--Бокса,
        Q--Q. Отчёт RMSE/MAPE на тесте, не только на обучении AIC.
\end{enumerate}
Типичный итог: SARIMA и HW близки по RMSE. SARIMA даёт более узкие интервалы при белых остатках. Значимый Q~--- повод упростить или расширить модель, а не подогнать $(p,q)$.

Семинар~10: SPRT Бернулли и сравнение с фиксированным $n$. Warmup. Задают коридор $p_L=0{,}45$, $p_U=0{,}55$, уровни $\alpha,\beta$, вычисляют $A=(1-\beta)/\alpha$, $B=\beta/(1-\alpha)$ и границы $a_m,r_m$ для SPRT по числу успехов $d_m=\sum X_i$.

Иллюзия фиксированного теста. Симулируют $X_i\sim\mathrm{Ber}(p)$ и на каждом $m$ считают p-value \texttt{proportions\_ztest} против $p_0=0{,}5$ без поправки за многократный просмотр.
При $p=0{,}4$ или $p=0{,}6$ подглядывание даёт ранние ложные отклонения $H_0$. При $p=0{,}5$ тест значим только после большого $m$.

SPRT (стиль вакцины / A/B конверсии). Библиотека \texttt{sprt}, \texttt{SPRTBinomial(alpha, beta, h0=0.45, h1=0.55, values=...)}.
Останавливаются, когда $d_m$ пересекает $a_m$ или $r_m$. Фиксируют фактическое $N$~--- случайная величина, её нужно сообщать в отчёте.

Нормальный SPRT (ADHD). Ряд времени реакции. Тестируют, уменьшает ли вмешательство среднее время относительно контроля, задавая $\mu_L,\mu_U$ как практически значимый эффект (\texttt{SPRTNormal} с известной или оценённой дисперсией). Автокорреляцию по времени нужно учесть отдельно (агрегация в блоки или модель остатков).

Сравнение последовательного и фиксированного-$n$. На семинаре сопоставляют:
\begin{itemize}
  \item фиксированный z-тест доли при повторных взглядах~--- завышенный
        уровень ошибки I рода;
  \item SPRT с заранее заданными $p_L,p_U$~--- контроль $\alpha,\beta$ при
        корректной остановке;
  \item двухэтапный и полностью последовательный ДИ Стейна для среднего~---
        фиксированная длина интервала ценой случайного $N$.
\end{itemize}
Содержательный вывод: в мониторинге экспериментов процедуру фиксируют до сбора данных. Иначе это частный случай множественной проверки (гл.~\ref{ch:multiple}).

\section{Контрольные вопросы по теме}

\begin{enumerate}
    \item Чем ACF на рис.~\ref{fig:ts-acf} отличается от корреляции
          в обычной регрессии?
    \item Запишите аддитивную и мультипликативную декомпозицию. Когда
          уместен $\log y_t$?
    \item Чем KPSS дополняет ADF?
    \item Запишите AR($p$), MA($q$), ARMA и SARIMA в операторной форме.
    \item Как устроен Q-критерий Льюнга--Бокса и почему $K$ вычитается
          из степеней свободы?
    \item Чем отличаются доверительный и предсказательный интервалы
          в ARIMA после преобразования $y$?
    \item Объясните смысл границ $a_m$ и $r_m$ и формулу $\mathbb{E}_p(N)$
          в SPRT для Бернулли.
    \item Зачем нужны усечение и batching? Чем отличается двухэтапный
          интервал Стейна от полностью последовательного?
    \item Назовите три типичные ошибки A/B-тестирования.
\end{enumerate}
